\section{RCOUNT Inputs And Report Contract}

The required inputs are RCOUNT reporting units, normalized election features,
and geography sufficient to define shared-boundary adjacency. The input package
also supplies source hashes, reporting-unit lineage, and any feature provenance
needed to reproduce the score. If a jurisdiction supplies geography at a
different level than results, the lineage and aggregation rule are part of the
feature set, not an informal preprocessing step.

W.05 instantiates the W.01 forensic report contract with spatial-outlier
semantics:

\begin{center}
\small
\begin{tabularx}{\textwidth}{lX}
\toprule
Field & W.05 meaning \\
\midrule
\texttt{method\_id} & Registered W.05 spatial-outlier analytic. \\
\texttt{feature\_set} & Unit features scored, such as turnout, vote share, residuals, or discrepancy rates. \\
\texttt{baseline\_population} & The neighbor set, plus any wider comparison population or time window. \\
\texttt{model\_id} & \texttt{spatial-outlier}, with the specific statistic or residual model named. \\
\texttt{unit\_id} & Reporting unit being scored. \\
\texttt{score} & Local outlier statistic, residual, or neighbor-difference score. \\
\texttt{p\_value\_or\_rank} & Optional calibration, permutation p-value, percentile, or queue rank. \\
\texttt{investigation\_note} & Machine-readable caveat, including edge effects or known boundary explanations. \\
\bottomrule
\end{tabularx}
\end{center}

The report metadata binds the output to source package hashes, geography hashes,
feature definitions, adjacency construction rules, and model parameters. That
binding is essential: a second analyst with the same RCOUNT package and the same
registered method should be able to recompute the same investigation queue. The
contract is defined in W.01; W.05 only specializes it for neighbor anomalies.
